\chapter{Algorithmes}\label{ch:iv3:algorithms}

\begin{ivtrains}
Choisir la structure de données avant d'écrire la boucle, en énoncer le coût et traiter correctement les cas limites~: entrée vide, égalités, doublons et valeurs pouvant être négatives. Les prix, les carnets d'ordres et les flux n'apparaissent que là où ils changent le problème. Chaque solution s'écrit comme on le ferait au tableau, en Python ou en C++, puis on en indique le coût.
\end{ivtrains}

\begin{question}[$\star$ \iqroles{QD, Dev, QR, QT} \iqfirm{toutes} \iqlevel{junior}]\label{iq:iv3:algorithms:1}
On vous donne une liste de prix classés dans l'ordre chronologique. Le drawdown à un instant donné est le prix le plus élevé observé jusque-là moins le prix actuel. Écrivez une fonction qui renvoie le plus grand drawdown de la liste, en un seul parcours et avec une mémoire supplémentaire constante. Que renvoie-t-elle pour \texttt{[100, 120, 90, 110, 80, 130]}~?
\end{question}

\begin{question}[$\star$ \iqroles{QD, Dev, LLE} \iqfirm{toutes} \iqlevel{junior}]\label{iq:iv3:algorithms:2}
Étant donné une liste d'entiers et une cible, renvoyez les positions de deux éléments distincts dont la somme vaut la cible, ou signalez qu'il n'en existe aucun. Elle doit s'exécuter en un temps proportionnel à la longueur de la liste en moyenne. Que doit-elle renvoyer pour la liste \texttt{[3, 4]} et la cible 6~?
\end{question}

\begin{question}[$\star$ \iqroles{QD, Dev, LLE, QT} \iqfirm{teneur de marché, crypto, hedge fund} \iqlevel{junior}]\label{iq:iv3:algorithms:3}
Le sell side d'un carnet d'ordres est une liste de niveaux de prix, chacun composé d'un prix et d'une taille, du prix le plus bas vers le haut. Un ordre d'achat au marché prend d'abord les actions les moins chères. Écrivez une fonction qui prend les niveaux et une quantité, et renvoie les exécutions, la quantité non exécutée et le coût total. Pour les niveaux \texttt{(100, 300), (101, 200), (103, 500)} et un ordre de 600 actions, quelles sont les exécutions et le prix moyen payé~?
\end{question}

\begin{question}[$\star$ \iqroles{QD, Dev, LLE, QR} \iqfirm{toutes} \iqlevel{junior}]\label{iq:iv3:algorithms:4}
Écrivez une classe dont la méthode \texttt{update(x)} prend le nombre suivant d'un flux et renvoie la moyenne des $k$ derniers nombres, ou de tous les nombres tant qu'il y en a moins de $k$. Chaque mise à jour doit prendre un temps constant, et la classe peut conserver une mémoire en $O(k)$. Déroulez ensuite la classe mentalement avec $k = 3$ sur les flottants Python \texttt{1e16, 1.0, 1.0, 1.0}~: que renvoie le dernier appel si la classe tient un total courant, en retranchant la valeur qui sort puis en ajoutant la nouvelle, et quelle est la vraie moyenne~?
\end{question}

\begin{question}[$\star$ \iqroles{QD, Dev, QR} \iqfirm{toutes} \iqlevel{junior}]\label{iq:iv3:algorithms:5}
Une chaîne de caractères ne contient que les caractères \texttt{(}, \texttt{)}, \texttt{[}, \texttt{]}, \texttt{\{} and \texttt{\}}. Écrivez une fonction qui décide si chaque caractère d'ouverture est refermé par le caractère correspondant, dans le bon ordre. Que renvoie-t-elle pour \texttt{([)]}, et pour la chaîne vide~?
\end{question}

\begin{question}[$\star$ \iqroles{QD, Dev, LLE} \iqfirm{toutes} \iqlevel{junior}]\label{iq:iv3:algorithms:6}
Une liste triée de $n$ prix peut contenir des répétitions. Écrivez une fonction qui renvoie la première position dont la valeur est au moins égale à $x$, ou $n$ s'il n'y en a aucune, en temps $O(\log n)$. Que renvoie-t-elle pour \texttt{[1, 2, 2, 2, 3]} avec $x = 2$, puis avec $x = 4$~? En C++, avec des positions de type \texttt{int}, que ne va-t-il pas avec le calcul du milieu \texttt{(lo + hi) / 2}~?
\end{question}

\begin{question}[$\star$ \iqroles{QD, Dev} \iqfirm{toutes} \iqlevel{junior}]\label{iq:iv3:algorithms:7}
Des symboles arrivent un à un. Écrivez une classe dotée de \texttt{add(symbol)} et \texttt{first\_unique()}, où la seconde méthode renvoie le premier symbole vu exactement une fois jusqu'ici, ou rien s'il n'y en a aucun. Les deux doivent prendre un temps constant en moyenne. Que renvoie-t-elle après la suite A, B, A, C, B~?
\end{question}

\begin{question}[$\star$ \iqroles{QD, Dev, LLE} \iqfirm{toutes} \iqlevel{junior}]\label{iq:iv3:algorithms:8}
Chaque nœud d'une liste chaînée simple porte une valeur et un pointeur vers le nœud suivant. Inversez la liste sur place, avec une mémoire supplémentaire constante, et renvoyez la nouvelle tête. Que doit-elle faire pour une liste vide~? En Python, que se passe-t-il pour la version récursive courte sur une liste de 100\,000 nœuds avec les réglages par défaut~?
\end{question}

\begin{question}[$\star$ \iqroles{LLE, QD, Dev} \iqfirm{toutes} \iqlevel{junior}]\label{iq:iv3:algorithms:9}
Pour un entier positif \texttt{x}, que calcule l'expression \texttt{x \& (x - 1)}~? Utilisez-la pour tester si \texttt{x} est une puissance de deux, et pour compter ses bits à 1. Arrondissez ensuite une taille demandée à la puissance de deux supérieure, et expliquez pourquoi un tampon circulaire (ring buffer) exige cette taille.
\end{question}

\begin{question}[$\star$ \iqroles{QD, Dev, LLE} \iqfirm{toutes} \iqlevel{junior}]\label{iq:iv3:algorithms:10}
Un fichier doit contenir chaque numéro de séquence de 1 à $n$, dans un ordre quelconque, mais exactement un manque. Trouvez-le en un seul parcours avec une mémoire supplémentaire constante. Que renvoyez-vous pour \texttt{[3, 1, 5, 2]} avec $n = 5$~? Si vous additionnez les nombres dans un entier signé de 64 bits, peut-il y avoir un dépassement de capacité~? Et si deux manquent~?
\end{question}

\begin{question}[$\star$ \iqroles{QD, Dev} \iqfirm{toutes} \iqlevel{junior}]\label{iq:iv3:algorithms:11}
Concevez une pile dotée de \texttt{push}, \texttt{pop} et \texttt{min}, chacune en temps constant dans le pire des cas. Les valeurs peuvent se répéter. Pourquoi une seule variable contenant le minimum courant ne suffit-elle pas~?
\end{question}

\begin{question}[$\star$ \iqroles{QD, Dev, LLE} \iqfirm{toutes} \iqlevel{junior}]\label{iq:iv3:algorithms:12}
Construisez une file premier entré, premier sorti en n'utilisant que deux piles. Une opération \texttt{pop} peut déplacer tous les éléments~; pourquoi le coût reste-t-il constant par opération sur une longue suite~? Utiliseriez-vous cette structure sur un chemin où chaque opération isolée doit respecter une limite de latence~?
\end{question}

\begin{question}[$\star$ \iqroles{QD, Dev, LLE, QT} \iqfirm{teneur de marché, crypto, hedge fund} \iqlevel{junior}]\label{iq:iv3:algorithms:13}
Un flux envoie des prix sous forme de texte décimal, comme 101.25, avec un tick de 0.01. Écrivez une fonction qui convertit ce texte en un nombre entier de ticks, sans virgule flottante, et qui refuse un texte comportant plus de deux décimales. En Python, que donne \texttt{int(float(s) * 100)} pour le texte 1.15~? Que doit renvoyer votre fonction pour le texte -0.05~?
\end{question}

\begin{question}[$\star$ \iqroles{QD, Dev, QR, QT} \iqfirm{toutes} \iqlevel{junior}]\label{iq:iv3:algorithms:14}
Vous disposez d'une liste de profits et pertes journaliers. Renvoyez la plus grande somme sur une plage non vide et contiguë de jours, ainsi que les jours où elle commence et où elle se termine, en un seul parcours. Que renvoyez-vous pour \texttt{[-2, 3, -1, 4, -5]}, et pour \texttt{[-3, -1, -2]}~?
\end{question}

\begin{question}[$\star$ \iqroles{QD, Dev, QR} \iqfirm{toutes} \iqlevel{junior}]\label{iq:iv3:algorithms:15}
Étant donné une liste d'entiers, renvoyez une liste dont la position $i$ contient le produit de tous les autres éléments. Faites-le en temps $O(n)$, sans division et sans tableau supplémentaire en dehors de la sortie. Que renvoyez-vous pour \texttt{[2, 3, 0, 4]}~? Pourquoi la méthode «~tout multiplier, puis diviser~» échoue-t-elle sur cette liste~?
\end{question}

\begin{question}[$\star$ \iqroles{QD, Dev, LLE, QT} \iqfirm{toutes} \iqlevel{junior}]\label{iq:iv3:algorithms:16}
Que donnent \texttt{-7 // 2} et \texttt{-7 \% 2} en Python 3, et \texttt{-7 / 2} et \texttt{-7 \% 2} pour un \texttt{int} en C++20~? Un calendar spread peut avoir un prix négatif. Écrivez une fonction, dans l'un ou l'autre langage, qui arrondit un prix entier vers le bas à un multiple d'un tick, de sorte que $-7$ avec un tick de 5 donne $-10$.
\end{question}

\begin{question}[$\star\star$ \iqroles{QD, Dev, LLE, QR} \iqfirm{toutes} \iqlevel{junior, confirmé}]\label{iq:iv3:algorithms:17}
Étant donné $n$ nombres et une taille de fenêtre $k$, renvoyez le plus grand nombre de chaque suite de $k$ nombres consécutifs. Le temps d'exécution doit être proportionnel à $n$, quelle que soit la valeur de $k$. Que renvoyez-vous pour \texttt{[1, 3, 2, 5, 4, 1]} avec $k = 3$~? Une étape de votre boucle peut effectuer beaucoup de travail~; comment justifiez-vous ce temps d'exécution~?
\end{question}

\begin{question}[$\star\star$ \iqroles{QD, Dev, QR} \iqfirm{toutes} \iqlevel{junior, confirmé}]\label{iq:iv3:algorithms:18}
Une liste d'entiers peut contenir des valeurs négatives. Comptez les plages contiguës de la liste dont la somme vaut exactement $T$, en temps proportionnel à la longueur de la liste en moyenne. Pour \texttt{[1, -1, 1]} et $T = 1$, quel est le compte~? Pourquoi une fenêtre à deux pointeurs, qui s'agrandit tant que la somme est inférieure à $T$ et se réduit tant qu'elle est supérieure, ne fonctionne-t-elle pas ici~?
\end{question}

\begin{question}[$\star\star$ \iqroles{QD, Dev, LLE} \iqfirm{toutes} \iqlevel{junior, confirmé}]\label{iq:iv3:algorithms:19}
Chaque ordre possède une heure d'ouverture entière et une heure de clôture entière ultérieure, et il est actif à l'instant $t$ lorsque $\text{ouverture} \le t < \text{clôture}$. Étant donné $n$ ordres, dans un ordre quelconque, trouvez le plus grand nombre d'ordres actifs à la fois et le premier instant où ce nombre est atteint, en temps $O(n \log n)$. Que renvoyez-vous pour les deux ordres $(1, 2)$ et $(2, 3)$~?
\end{question}

\begin{question}[$\star\star$ \iqroles{QD, Dev, QR, DS} \iqfirm{toutes} \iqlevel{junior, confirmé}]\label{iq:iv3:algorithms:20}
Vous disposez de $n$ transactions, chacune avec un symbole, un prix et une quantité, et $m$ symboles distincts. Renvoyez les $k$ symboles au notionnel total le plus élevé (prix multiplié par quantité), du plus grand au plus petit, en départageant les égalités par nom de symbole. Donnez une solution qui coûte moins que trier les $m$ symboles lorsque $k$ est petit. Que se passe-t-il si les transactions continuent d'arriver et qu'il faut répondre à tout instant~?
\end{question}

\begin{question}[$\star\star$ \iqroles{QD, Dev, LLE} \iqfirm{teneur de marché, crypto, hedge fund} \iqlevel{junior, confirmé}]\label{iq:iv3:algorithms:21}
En Python, vous disposez de $K$ flux, chacun étant un itérateur de couples \texttt{(timestamp, message)} dans l'ordre temporel, où chaque message est un dictionnaire. Écrivez un générateur qui produit tous les couples dans l'ordre temporel. À horodatages égaux, le flux de numéro le plus bas passe en premier, et le générateur ne peut garder en mémoire qu'un seul couple par flux. Quel est son coût pour $N$ couples au total~? Que devient votre générateur lorsqu'un flux reste silencieux pendant une heure~?
\end{question}

\begin{question}[$\star\star$ \iqroles{QD, Dev, QR, DS} \iqfirm{toutes} \iqlevel{junior, confirmé}]\label{iq:iv3:algorithms:22}
Des nombres arrivent un à un. Concevez une classe dotée de \texttt{add(x)} et \texttt{median()}, où \texttt{median} renvoie la valeur médiane, ou la moyenne des deux valeurs médianes lorsque le nombre est pair. Indiquez le coût de chaque méthode. Que changeriez-vous si chaque valeur était un nombre entier de ticks compris entre 0 et 1000~?
\end{question}

\begin{question}[$\star\star$ \iqroles{QD, Dev, QR, DS} \iqfirm{toutes} \iqlevel{junior, confirmé}]\label{iq:iv3:algorithms:23}
Un flux de longueur inconnue ne passe qu'une seule fois. Conservez exactement $k$ éléments en mémoire, de sorte que, lorsque le flux se termine, chaque sous-ensemble de $k$ éléments parmi ceux qui sont passés ait la même probabilité d'être celui que vous gardez. Écrivez la fonction et démontrez-la, au moins pour la probabilité qu'un élément donné soit conservé.
\end{question}

\begin{question}[$\star\star$ \iqroles{QD, Dev, LLE} \iqfirm{toutes} \iqlevel{junior, confirmé}]\label{iq:iv3:algorithms:24}
En C++20, écrivez un cache à capacité fixe offrant deux opérations~: \texttt{get(key)} renvoie la valeur si elle est présente, et \texttt{put(key, value)} insère ou écrase. Lorsqu'il est plein, l'insertion d'une nouvelle clé élimine la clé dont la dernière utilisation est la plus ancienne, \texttt{get} et \texttt{put} comptant tous deux comme une utilisation. Les deux opérations doivent prendre un temps constant en moyenne. Pourquoi \texttt{std::vector} ou \texttt{std::deque} ne conviendraient-ils pas pour l'ordre de récence~?
\end{question}

\begin{question}[$\star\star$ \iqroles{QD, Dev, LLE} \iqfirm{toutes} \iqlevel{junior}]\label{iq:iv3:algorithms:25}
Une liste chaînée simple peut se terminer par un cycle, de sorte que suivre les pointeurs n'atteint jamais de fin. Avec une mémoire supplémentaire constante, décidez si c'est le cas et renvoyez le nœud où le cycle commence. Expliquez pourquoi votre seconde phase tombe exactement sur ce nœud.
\end{question}

\begin{question}[$\star\star$ \iqroles{QD, Dev, QR} \iqfirm{toutes} \iqlevel{junior}]\label{iq:iv3:algorithms:26}
Vous disposez d'une liste de suspensions de cotation, chacune avec une heure de début et une heure de fin, dans un ordre quelconque~; deux suspensions qui se touchent ou se chevauchent ne comptent que pour une. Renvoyez les suspensions fusionnées dans l'ordre temporel, ainsi que la durée totale de suspension. Que renvoyez-vous pour $(8, 10)$, $(1, 3)$, $(10, 12)$, $(2, 6)$~? Et si les suspensions arrivent une à une et qu'il faut la liste fusionnée après chacune~?
\end{question}

\begin{question}[$\star\star$ \iqroles{QD, Dev, QR} \iqfirm{toutes} \iqlevel{junior}]\label{iq:iv3:algorithms:27}
Étant donné une liste de cours quotidiens, renvoyez pour chaque jour le nombre de jours jusqu'à un cours strictement supérieur, ou 0 s'il n'en vient aucun. Le calcul doit s'exécuter en temps $O(n)$. Que renvoyez-vous pour \texttt{[100, 101, 103, 99, 97, 100, 104, 101]}~? Une étape peut faire beaucoup de travail~; pourquoi le total reste-t-il linéaire~?
\end{question}

\begin{question}[$\star\star$ \iqroles{QD, Dev, QT} \iqfirm{toutes} \iqlevel{junior}]\label{iq:iv3:algorithms:28}
Vous devez envoyer exactement $n$ actions, et chaque ordre enfant doit avoir une taille prise dans une liste fixe. Trouvez le nombre minimal d'ordres enfants et leurs tailles, ou signalez qu'il est impossible de le faire. Avec les tailles 1, 3 et 4 et $n = 6$, que donne la méthode qui prend à chaque fois la plus grande taille qui tient, et quelle est la vraie réponse~? Dites quel est le coût de votre méthode.
\end{question}

\begin{question}[$\star\star$ \iqroles{QD, Dev, QR, QT} \iqfirm{toutes} \iqlevel{junior, confirmé}]\label{iq:iv3:algorithms:29}
Vous disposez d'une liste de cours quotidiens et ne pouvez détenir qu'une action à la fois. Vous pouvez acheter et vendre aussi souvent que vous le souhaitez, et chaque aller-retour complet entraîne des frais $f$. Renvoyez le plus grand profit total en temps $O(n)$ et en mémoire constante. Quel résultat obtenez-vous pour \texttt{[10, 13, 12, 18, 14, 19]} avec $f = 2$~? Que renvoie votre méthode sans frais, et avec des frais d'un million~?
\end{question}

\begin{question}[$\star\star$ \iqroles{QD, Dev, DS} \iqfirm{toutes} \iqlevel{junior}]\label{iq:iv3:algorithms:30}
Vous rapprochez deux listes de noms de titres financiers. Écrivez une fonction calculant la distance d'édition entre deux chaînes~: le nombre minimal d'insertions, de suppressions et de substitutions d'un seul caractère permettant de transformer l'une en l'autre. Elle doit s'exécuter en temps $O(mn)$ et n'utiliser que $O(\min(m, n))$ de mémoire. Quelle est la distance de ABCD à ACBD~? Que perd-on en ne gardant qu'une seule ligne~?
\end{question}

\begin{question}[$\star\star$ \iqroles{QD, Dev, QR} \iqfirm{toutes} \iqlevel{junior, confirmé}]\label{iq:iv3:algorithms:31}
Il y a $n$ comptes. Des paires arrivent une à une, chacune indiquant que deux comptes appartiennent au même groupe, et les liens se transmettent~: si $a$ appartient au même groupe que $b$ et $b$ au même groupe que $c$, alors $a$ appartient au même groupe que $c$. À tout moment, vous devez répondre si deux comptes sont dans un même groupe et combien il y a de groupes. Écrivez-le, et dites quel est son coût par opération. Que ne peut pas faire votre structure~?
\end{question}

\begin{question}[$\star\star$ \iqroles{QD, Dev, QR} \iqfirm{toutes} \iqlevel{junior, confirmé}]\label{iq:iv3:algorithms:32}
Un lot nocturne comporte $n$ tâches, et une liste de paires $(a, b)$ signifiant que $a$ doit se terminer avant que $b$ ne commence. Renvoyez un ordre qui les respecte toutes, ou signalez qu'aucun n'existe, en temps linéaire en nombre de tâches et de paires. Donnez maintenant une durée à chaque tâche, un nombre quelconque de tâches pouvant s'exécuter en même temps. Les tâches A, B, C et D durent respectivement 3, 2, 4 et 1 heures, B et C dépendent de A, et D dépend de B et C. Quand tout est-il terminé~?
\end{question}

\begin{question}[$\star\star$ \iqroles{QD, Dev} \iqfirm{toutes} \iqlevel{junior}]\label{iq:iv3:algorithms:33}
Une grille de $r$ lignes et $c$ colonnes comporte des cases barrées par des murs. Vous vous déplacez d'une case à la fois, vers le haut, le bas, la gauche ou la droite. Trouvez le nombre minimal de déplacements d'une case de départ à une case d'arrivée, ou signalez que cette dernière est inatteignable. Pour les lignes \texttt{.{}.\#}, \texttt{.\#.} et \texttt{.{}.{}.} (avec \texttt{\#} désignant un mur), en allant du coin supérieur gauche au coin inférieur droit, quelle est la réponse~? Pourquoi le parcours en profondeur ne donne-t-il pas ce résultat~? Importe-t-il de marquer une case quand on l'ajoute ou quand on la retire~?
\end{question}

\begin{question}[$\star\star$ \iqroles{QD, Dev, LLE} \iqfirm{toutes} \iqlevel{junior, confirmé}]\label{iq:iv3:algorithms:34}
Des plateformes de négociation sont reliées par des liens unidirectionnels, chacun ayant une latence en microsecondes jamais négative. Trouvez l'itinéraire le plus rapide depuis une plateforme vers chacune des autres, avec l'itinéraire lui-même, et dites quel est son coût pour $V$ plateformes et $E$ liens. Avec les liens A vers B de 5, A vers C de 2, C vers B de 1 et B vers D de 1, quel est le temps minimal de A à D, et par quel itinéraire~? Qu'est-ce qui ne va pas si un lien a un poids négatif~?
\end{question}

\begin{question}[$\star\star$ \iqroles{QD, Dev, QR} \iqfirm{toutes} \iqlevel{junior, confirmé}]\label{iq:iv3:algorithms:35}
Une liste de tranches d'ordres, toutes de taille positive, doit être envoyée dans l'ordre donné sur $d$ jours, chaque jour prenant une suite contiguë de tranches. Minimisez le total du jour le plus chargé. Quelle est la réponse pour \texttt{[4, 9, 3, 7, 8]} avec $d = 2$, et avec $d = 3$~? Trouvez une méthode plus rapide que d'essayer tous les découpages, et dites pourquoi elle exige que les tailles soient positives.
\end{question}

\begin{question}[$\star\star$ \iqroles{QD, Dev, QR} \iqfirm{toutes} \iqlevel{junior, confirmé}]\label{iq:iv3:algorithms:36}
Trouvez le $k$-ième plus petit parmi $n$ nombres en temps moyen $O(n)$, sans les trier tous. Pourquoi votre méthode devient-elle quadratique sur une liste triée si le pivot est toujours le premier élément~? Que se passe-t-il sur une liste où toutes les valeurs sont égales, si l'on partitionne en «~inférieur au pivot~» et «~non inférieur au pivot~»~? Comment évitez-vous les deux~?
\end{question}

\begin{question}[$\star\star$ \iqroles{QD, Dev} \iqfirm{toutes} \iqlevel{junior}]\label{iq:iv3:algorithms:37}
Étant donné une liste d'entiers, renvoyez chaque triplet de valeurs distinctes dont la somme vaut zéro, chacun une seule fois, en temps $O(n^2)$. Que renvoyez-vous pour \texttt{[-2, 0, 1, 1, 2, -1, -1]}~? D'où viennent les doublons, et comment les écarter sans ensemble de résultats~?
\end{question}

\begin{question}[$\star\star$ \iqroles{QD, Dev, DS} \iqfirm{toutes} \iqlevel{junior}]\label{iq:iv3:algorithms:38}
Le tape est une liste de symboles, un par transaction. Renvoyez la longueur de la plus longue portion contiguë contenant au plus $k$ symboles distincts, en temps $O(n)$. Que renvoyez-vous pour le tape A, B, A, C, C, D, D, D avec $k = 2$~? Comment compteriez-vous les portions contenant exactement $k$ symboles distincts~?
\end{question}

\begin{question}[$\star\star$ \iqroles{LLE, QD, Dev} \iqfirm{teneur de marché, hedge fund, crypto} \iqlevel{junior, confirmé}]\label{iq:iv3:algorithms:39}
En C++20, écrivez une table de hachage associant des clés non nulles de 64 bits à des valeurs de 64 bits, avec sondage linéaire dans un tableau unique et plat dont la taille est une puissance de deux, qui n'alloue rien après sa construction. Elle possède \texttt{put}, \texttt{get} et \texttt{erase}. Pourquoi le simple fait de «~vider la case~» constitue-t-il un mauvais \texttt{erase}, et que faites-vous à la place~? Pourquoi la table doit-elle rester loin d'être pleine~?
\end{question}

\begin{question}[$\star\star$ \iqroles{QD, Dev, LLE} \iqfirm{toutes} \iqlevel{junior}]\label{iq:iv3:algorithms:40}
Implémentez un tas binaire minimal dans un tableau, avec \texttt{push}, \texttt{pop} et un constructeur à partir de $n$ éléments. Dites ce que coûte chacune de ces opérations. Pourquoi construire le tas par des descentes depuis le dernier parent coûte-t-il $O(n)$, alors que $n$ ajouts successifs peuvent coûter $O(n \log n)$~?
\end{question}

\begin{question}[$\star\star$ \iqroles{QD, Dev, LLE} \iqfirm{toutes} \iqlevel{junior, confirmé}]\label{iq:iv3:algorithms:41}
Dans une liste d'identifiants d'ordres, chaque identifiant apparaît exactement deux fois, sauf deux identifiants qui n'apparaissent qu'une fois. Trouvez ces deux identifiants en temps $O(n)$ et en mémoire auxiliaire constante. Que renvoyez-vous pour \texttt{[4, 7, 4, 9, 7, 2]}~? Et si chaque identifiant apparaît trois fois, sauf un~?
\end{question}

\begin{question}[$\star\star$ \iqroles{QD, Dev, LLE} \iqfirm{teneur de marché, crypto, banque} \iqlevel{junior, confirmé}]\label{iq:iv3:algorithms:42}
Une bourse accepte au plus $N$ messages dans toute fenêtre de $T$ secondes. Écrivez \texttt{allow(now)}, avec des temps qui ne décroissent jamais, de sorte que vous ne dépassiez jamais la limite. Avec $N = 2$ et $T = 10$, et une fenêtre qui conserve les temps supérieurs à \texttt{now - T} jusqu'à \texttt{now}, que répondez-vous pour les temps 0, 1, 10, 10 et 11~? Pourquoi un compteur remis à zéro toutes les $T$ secondes ne suffit-il pas, et comment se compare un seau à jetons (token bucket)~?
\end{question}

\begin{question}[$\star\star$ \iqroles{QD, Dev, DS} \iqfirm{toutes} \iqlevel{junior}]\label{iq:iv3:algorithms:43}
Écrivez une structure de stockage dotée de \texttt{set(key, value, t)} et \texttt{get(key, t)}, où \texttt{get} renvoie la valeur dont le temps est le plus récent parmi ceux antérieurs ou égaux à $t$, ou aucune valeur. Pour chaque clé, les temps arrivent par ordre croissant. Quel est le coût de chaque appel~? Qu'est-ce qui change si les temps peuvent arriver dans le désordre~? Comment joindriez-vous deux séries temporelles triées de cette façon, en temps linéaire~?
\end{question}

\begin{question}[$\star\star$ \iqroles{QD, Dev} \iqfirm{toutes} \iqlevel{junior}]\label{iq:iv3:algorithms:44}
Concevez un ensemble d'identifiants d'ordres muni de \texttt{insert}, \texttt{remove} et \texttt{pick}, où \texttt{pick} renvoie un identifiant tiré uniformément au hasard. Les trois opérations doivent prendre un temps constant en moyenne. Pourquoi une table de hachage seule ne suffit-elle pas~? Quel est le bogue classique dans \texttt{remove}~?
\end{question}

\begin{question}[$\star\star$ \iqroles{QD, Dev, QR} \iqfirm{toutes} \iqlevel{junior}]\label{iq:iv3:algorithms:45}
Écrivez un mélange de $n$ éléments qui rend toutes les $n!$ permutations équiprobables, sur place et en temps $O(n)$, à partir d'une fonction d'entiers aléatoires uniformes. Un collègue procède autrement~: il échange chaque position avec une position choisie uniformément dans toute la liste. Pour trois éléments, combien de résultats aléatoires équiprobables cela donne-t-il, et pourquoi cela ne peut-il pas être équitable~?
\end{question}

\begin{question}[$\star\star$ \iqroles{QD, Dev} \iqfirm{toutes} \iqlevel{junior}]\label{iq:iv3:algorithms:46}
Un arbre binaire contient des clés distinctes. Décidez si c'est un arbre binaire de recherche valide~: tout ce qui se trouve dans le sous-arbre gauche de chaque nœud est inférieur à ce nœud, et tout ce qui se trouve dans son sous-arbre droit lui est supérieur. Pourquoi comparer chaque nœud seulement à ses deux enfants échoue-t-il~? Donnez un arbre qui passe ce test et qui n'est pas valide. Et qu'en est-il d'un arbre de 100\,000 niveaux de profondeur~?
\end{question}

\begin{question}[$\star\star$ \iqroles{QR, QD, DS} \iqfirm{hedge fund, teneur de marché, banque} \iqlevel{junior, confirmé}]\label{iq:iv3:algorithms:47}
Écrivez une classe qui prend un flux de flottants et renvoie la moyenne courante et la variance d'échantillon, en temps et en mémoire constants par valeur. Pourquoi ne pas simplement garder le compte, la somme et la somme des carrés~? Qu'est-ce que cela donne en double précision pour les quatre valeurs $10^9 + 4$, $10^9 + 7$, $10^9 + 13$ et $10^9 + 16$, dont la vraie variance d'échantillon vaut 30~? Comment traitez-vous une version à pondération exponentielle~?
\end{question}

\begin{question}[$\star\star$ \iqroles{LLE, QD, Dev} \iqfirm{teneur de marché, crypto, hedge fund} \iqlevel{junior, confirmé}]\label{iq:iv3:algorithms:48}
Une colonne de prix entiers évolue par petits pas. Stockez-la sous forme de différences, chaque différence étant transformée en entier non signé par la règle zigzag (0, $-1$, 1, $-2$ deviennent 0, 1, 2, 3), puis écrite sous forme de varint (sept bits par octet, le bit de poids fort étant à 1 sur tous les octets sauf le dernier, les bits de poids faible en premier). Écrivez l'encodeur et le décodeur. Combien d'octets occupe \texttt{[1000, 1001, 999, 1300]}~? Pourquoi le zigzag, au juste~?
\end{question}

\begin{question}[$\star\star\star$ \iqroles{QD, Dev, LLE} \iqfirm{teneur de marché, crypto, hedge fund} \iqlevel{junior, confirmé}]\label{iq:iv3:algorithms:49}
Un marché reçoit des ordres d'achat et de vente, chacun avec un identifiant, un prix et une quantité. En C++20, écrivez une structure capable d'ajouter un ordre, d'annuler un ordre par son identifiant, et de renvoyer le meilleur cours acheteur et le meilleur cours vendeur, avec la quantité totale à chacun. Indiquez ce que coûte chaque opération, $L$ étant le nombre de niveaux de prix. Les identifiants en double et les annulations d'identifiants inconnus doivent être refusés. Il n'y a pas d'appariement~: les deux côtés restent simplement là.
\end{question}

\begin{question}[$\star\star\star$ \iqroles{QD, Dev, QR, QT} \iqfirm{banque, crypto, teneur de marché} \iqlevel{junior, confirmé}]\label{iq:iv3:algorithms:50}
On vous donne un tableau de taux de change~: pour certaines paires ordonnées de devises $(a, b)$, une unité de $a$ achète $r_{ab}$ unités de $b$. Un arbitrage est une suite de conversions qui commence et se termine dans la même devise avec davantage qu'au départ, et qui peut comporter un nombre quelconque d'étapes. Écrivez une fonction qui renvoie une telle suite si elle existe, et indiquez son coût pour $n$ devises. Pourquoi vérifier chaque triangle ne suffit-il pas~?
\end{question}

\begin{question}[$\star\star\star$ \iqroles{QD, Dev, QR} \iqfirm{toutes} \iqlevel{junior, confirmé}]\label{iq:iv3:algorithms:51}
Renvoyez une plus longue sous-suite strictement croissante d'une liste de $n$ nombres~: non pas sa seule longueur, mais les éléments eux-mêmes, en temps $O(n \log n)$. Les éléments n'ont pas besoin d'être adjacents. Essayez votre méthode sur \texttt{[3, 4, 1]}.
\end{question}

\begin{question}[$\star\star\star$ \iqroles{QR, QD, DS} \iqfirm{hedge fund, banque} \iqlevel{junior, confirmé}]\label{iq:iv3:algorithms:52}
Vous disposez de deux listes de $n$ nombres, appariées par position, toutes les valeurs de chaque liste étant distinctes. Deux positions sont concordantes si les deux listes les ordonnent de la même manière, et discordantes sinon. Le tau de Kendall est le nombre de paires concordantes moins le nombre de paires discordantes, divisé par le nombre de paires. Calculez-le en temps $O(n \log n)$, pour $n$ de l'ordre du million. Que faites-vous lorsque des valeurs peuvent être ex aequo~?
\end{question}

\begin{question}[$\star\star\star$ \iqroles{QD, Dev, LLE, DS} \iqfirm{teneur de marché, crypto, hedge fund} \iqlevel{confirmé}]\label{iq:iv3:algorithms:53}
Un flux de $n$ symboles passe une seule fois, $n$ n'étant pas connu à l'avance, et vous ne pouvez garder que $k - 1$ compteurs. Produisez un ensemble d'au plus $k - 1$ symboles qui contient à coup sûr chaque symbole apparaissant plus de $n / k$ fois. Écrivez-le, dites quel est son coût par élément, et démontrez la garantie. De combien vos compteurs peuvent-ils se tromper~? Pour $k = 2$, de quel problème bien connu s'agit-il~?
\end{question}

\begin{question}[$\star\star\star$ \iqroles{QD, Dev, QR, QT} \iqfirm{toutes} \iqlevel{junior, confirmé}]\label{iq:iv3:algorithms:54}
Vous disposez d'une liste de $n$ cours quotidiens et pouvez faire au plus $k$ allers-retours, en ne détenant qu'une action à la fois. Renvoyez le plus grand profit total en temps $O(nk)$. Quel résultat obtenez-vous pour \texttt{[4, 2, 7, 6, 1, 5]} avec $k = 1$ et avec $k = 2$~? Que devient-il lorsque $k$ est très grand, et dans quel cas un $k$ grand permet-il de se passer entièrement du tableau~?
\end{question}

\begin{question}[$\star\star\star$ \iqroles{QD, Dev} \iqfirm{toutes} \iqlevel{junior, confirmé}]\label{iq:iv3:algorithms:55}
Un histogramme est formé de barres de largeur 1 et de hauteurs données. Renvoyez la plus grande aire d'un rectangle contenu dans cet histogramme, en temps $O(n)$. Quel résultat obtenez-vous pour \texttt{[2, 1, 4, 5, 1, 3, 3]}~? Dites comment sont traitées les hauteurs égales, et pourquoi cela ne fait perdre aucun rectangle.
\end{question}

\begin{question}[$\star\star\star$ \iqroles{QD, Dev, LLE} \iqfirm{toutes} \iqlevel{confirmé}]\label{iq:iv3:algorithms:56}
Écrivez un cache LFU (least-frequently-used) de capacité fixe, dont les opérations \texttt{get} et \texttt{put} s'effectuent toutes deux en temps constant en moyenne. Chaque \texttt{get} ou \texttt{put} sur une clé compte comme une utilisation. Lorsqu'il est plein, une nouvelle clé évince la clé la moins utilisée et, parmi celles-ci, celle dont la dernière utilisation est la plus ancienne. Avec une capacité de 2, que reste-t-il après \texttt{put(a)}, \texttt{put(b)}, \texttt{get(a)}, \texttt{put(c)}~? Pourquoi une nouvelle clé remet-elle à 1 votre enregistrement du plus petit compte~?
\end{question}

\begin{question}[$\star\star\star$ \iqroles{QD, Dev, LLE, QT} \iqfirm{teneur de marché, crypto, hedge fund} \iqlevel{junior, confirmé}]\label{iq:iv3:algorithms:57}
En Python, écrivez un moteur d'appariement à ordres à cours limité, avec priorité prix-temps~: \texttt{submit(id, side, price, quantity)} exécute l'ordre contre le côté opposé au prix de l'ordre au repos, au meilleur prix d'abord et au plus ancien d'abord à prix égal, et tout reliquat se place au repos. \texttt{cancel(id)} retire un ordre au repos. Les identifiants sont uniques. Deux ventes sont au repos à 10, d'abord 100 actions puis 50 actions, et un achat de 120 à 11 arrive. Quelles transactions sont affichées, et que reste-t-il~? Que coûtent les deux appels~?
\end{question}

\begin{question}[$\star\star\star$ \iqroles{QD, Dev, QR} \iqfirm{toutes} \iqlevel{junior, confirmé}]\label{iq:iv3:algorithms:58}
Chaque tâche a un début, une fin et un profit, et deux tâches ne peuvent être retenues ensemble que si elles ne se chevauchent pas, une tâche pouvant commencer exactement au moment où une autre se termine. Renvoyez le plus grand profit total en temps $O(n \log n)$. Quel résultat obtenez-vous pour $(0, 3, 4)$, $(2, 5, 7)$, $(4, 7, 3)$ et $(6, 9, 5)$~? Pourquoi prendre d'abord la tâche la plus rentable échoue-t-il~?
\end{question}

\begin{question}[$\star\star\star$ \iqroles{QD, Dev, LLE} \iqfirm{teneur de marché, crypto, hedge fund} \iqlevel{junior, confirmé}]\label{iq:iv3:algorithms:59}
Le côté acheteur d'un carnet d'ordres est un tableau de $P$ niveaux de prix entiers, avec une taille à chacun. Il vous faut \texttt{update(price, delta)}, ainsi que \texttt{depth(price)}, la taille totale à ce prix ou mieux (c'est-à-dire à ce prix ou au-dessus), les deux en $O(\log P)$. Écrivez-le. Pourquoi un tableau simple échoue-t-il, et pourquoi un tableau de sommes cumulées échoue-t-il~? Comment trouveriez-vous le prix le plus élevé auquel la profondeur atteint une quantité donnée~?
\end{question}

\begin{question}[$\star\star\star$ \iqroles{QD, Dev, LLE} \iqfirm{toutes} \iqlevel{confirmé}]\label{iq:iv3:algorithms:60}
Les éléments d'un flux arrivent un à un. Étant donné un motif de $m$ éléments, signalez chaque instant où les $m$ derniers éléments égalent le motif, avec une mémoire en $O(m)$ et un temps constant amorti par élément. Pour le motif A, B, A, B et le flux A, B, A, B, A, B, à quelles positions, comptées à partir de 1, signalez-vous~? Montrez pourquoi revenir au début du motif après chaque non-correspondance est faux, et pas seulement lent.
\end{question}

\begin{question}[$\star\star\star$ \iqroles{QD, Dev, DS} \iqfirm{toutes} \iqlevel{junior, confirmé}]\label{iq:iv3:algorithms:61}
Une bande est une liste de symboles. Étant donné un multi-ensemble de symboles requis, qui peuvent se répéter, renvoyez le plus court segment contigu qui les contient tous, en temps $O(n)$. Pour la bande X, A, Y, B, Z, A, C, B et les symboles requis A, B, C, quel est-il~? Qu'est-ce qui change si A est requis deux fois~? Quel compteur vous indique que la fenêtre est devenue valide~?
\end{question}

\begin{question}[$\star\star\star$ \iqroles{LLE, QD, Dev, DS} \iqfirm{teneur de marché, crypto, hedge fund} \iqlevel{confirmé}]\label{iq:iv3:algorithms:62}
Vous voulez savoir si un identifiant d'ordre a déjà été vu, parmi 100 millions, avec moins de mémoire qu'une table de hachage. Concevez un filtre de Bloom~: un tableau de $m$ bits, où chaque clé active $h$ bits choisis par hachage, et une recherche teste ces bits. Combien de bits et combien de fonctions de hachage faut-il pour un taux de faux positifs de 1~\%~? Dérivez la formule. Quelle sorte d'erreur est impossible, et pourquoi ne peut-on pas supprimer~?
\end{question}

\begin{question}[$\star\star\star$ \iqroles{QD, Dev, LLE} \iqfirm{teneur de marché, crypto, hedge fund} \iqlevel{junior, confirmé}]\label{iq:iv3:algorithms:63}
Il existe $V$ plateformes, chacune avec son propre meilleur cours acheteur, qui peut monter ou descendre à tout moment. Écrivez \texttt{update(venue, price)} en $O(\log V)$, et \texttt{best()}, qui renvoie le cours acheteur le plus élevé et sa plateforme, en $O(1)$. Pourquoi le tas de la bibliothèque standard ne vous donne-t-il pas cela directement~? Quand $V = 8$, l'utiliseriez-vous vraiment~?
\end{question}

\begin{question}[$\star\star\star$ \iqroles{QD, Dev, DS} \iqfirm{toutes} \iqlevel{confirmé}]\label{iq:iv3:algorithms:64}
Vous devez trier un fichier de 10~To d'enregistrements de taille fixe, sur une machine disposant de 32~Go de mémoire et d'un seul disque qui lit ou écrit 500~Mo par seconde, en accès séquentiel. Décrivez la méthode. Combien de séries triées produisez-vous dans la première phase, combien de passes la fusion nécessite-t-elle, et combien de temps dure l'ensemble du travail si rien d'autre ne le limite~? Écrivez une petite version en Python.
\end{question}

\begin{question}[$\star\star\star$ \iqroles{QD, Dev, QR} \iqfirm{toutes} \iqlevel{confirmé}]\label{iq:iv3:algorithms:65}
Vous devez tirer des éléments d'une liste fixe de $n$ éléments de poids donnés, des millions de fois. La recherche dichotomique sur la somme cumulée coûte $O(\log n)$ par tirage. Montrez une méthode qui coûte $O(n)$ pour la préparation et $O(1)$ par tirage, en utilisant une position uniforme et un lancer de pile ou face. Construisez sa table pour les poids $\frac{1}{2}$, $\frac{1}{3}$ et $\frac{1}{6}$, et montrez qu'elle donne les bonnes probabilités.
\end{question}

\begin{question}[$\star\star\star$ \iqroles{QD, Dev, LLE, QT} \iqfirm{teneur de marché, crypto, hedge fund} \iqlevel{confirmé}]\label{iq:iv3:algorithms:66}
Un niveau de prix contient une file premier entré, premier sorti d'ordres. Les ordres rejoignent l'arrière de la file, s'annulent de n'importe où, et les transactions prélèvent des actions à l'avant. Il vous faut aussi dire, pour un ordre quelconque, combien d'actions le précèdent. Donnez une structure qui réalise les quatre opérations en $O(\log n)$. Après l'arrivée d'ordres de 100, 200 et 300 actions, l'annulation du deuxième et une transaction de 150 actions, combien d'actions précèdent le troisième~? Pourquoi ne pas conserver un total cumulé sur chaque ordre~?
\end{question}

\begin{question}[$\star\star\star$ \iqroles{QD, Dev, LLE} \iqfirm{teneur de marché, crypto, hedge fund} \iqlevel{confirmé}]\label{iq:iv3:algorithms:67}
Un feed envoie des mises à jour de carnet d'ordres sur un canal non fiable, chacune portant un numéro de séquence et fixant la taille d'un niveau de prix. Un second canal envoie de temps à autre un instantané de tout le carnet, étiqueté avec le dernier numéro de séquence qu'il inclut. Écrivez la logique qui maintient un carnet correct~: la prise de connexion en cours de route, la détection d'un trou de séquence et la reprise. Que faites-vous exactement d'une mise à jour qui arrive pendant que vous attendez un instantané~? Quelles mises à jour faut-il écarter, et lesquelles faut-il appliquer~?
\end{question}

\begin{question}[$\star\star\star$ \iqroles{QD, Dev, DS} \iqfirm{toutes} \iqlevel{junior, confirmé}]\label{iq:iv3:algorithms:68}
Vous disposez de séries triées de tailles 2, 3, 4 et 10. Fusionner deux séries coûte la somme de leurs tailles, et la série fusionnée prend leur place. Quel est le coût total minimal pour tout fusionner en une seule série, et combien coûte une fusion qui procède toujours à partir de la plus grande~? Écrivez la méthode et démontrez-la. Que faut-il ajouter si vous fusionnez trois séries à la fois~?
\end{question}

\begin{question}[$\star\star\star\star$ \iqroles{QD, Dev, QR} \iqfirm{teneur de marché, hedge fund} \iqlevel{confirmé}]\label{iq:iv3:algorithms:69}
Étant donné $n$ nombres et une taille de fenêtre $k$, renvoyez la médiane de chaque fenêtre de $k$ nombres consécutifs, soit comme valeur du milieu, soit comme moyenne des deux valeurs du milieu. Les nombres peuvent se répéter. Faites mieux que trier chaque fenêtre, et codez-la. Quels sont le temps et la mémoire de votre solution~? Dépendent-ils de $k$ ou de $n$~?
\end{question}

\begin{question}[$\star\star\star\star$ \iqroles{QD, Dev, QR} \iqfirm{toutes} \iqlevel{confirmé}]\label{iq:iv3:algorithms:70}
Une liste contient $n$ entiers, pouvant être négatifs, et $T$ est un entier positif. Renvoyez la longueur du plus court segment contigu dont la somme vaut au moins $T$, ou signalez qu'il n'en existe aucun, en temps $O(n)$. Pour \texttt{[84, -37, 32, 40, 95]} et $T = 167$, quelle est la réponse~? Montrez pourquoi votre méthode n'écarte jamais le bon début d'un segment.
\end{question}

\begin{question}[$\star\star\star\star$ \iqroles{QD, Dev, QR} \iqfirm{toutes} \iqlevel{confirmé}]\label{iq:iv3:algorithms:71}
On vous donne deux listes triées de longueurs $m$ et $n$, non toutes deux vides, qui peuvent avoir des longueurs très différentes et partager des valeurs. Renvoyez la médiane de l'ensemble de leurs valeurs en temps $O(\log \min(m, n))$, soit comme valeur du milieu, soit comme moyenne des deux valeurs du milieu. Supposez maintenant qu'il existe de nombreuses listes triées, sur des machines différentes, et un billion d'entiers compris entre 0 et $10^9$~: comment trouvez-vous la médiane~?
\end{question}

\begin{question}[$\star\star\star\star$ \iqroles{QD, Dev, LLE} \iqfirm{teneur de marché, crypto, hedge fund} \iqlevel{confirmé}]\label{iq:iv3:algorithms:72}
Chaque cours acheteur a une heure d'ouverture, une heure de fermeture ultérieure et un prix, et il figure dans le carnet sur l'intervalle semi-ouvert allant de l'ouverture à la fermeture. Étant donné $n$ cours acheteurs dans un ordre quelconque, produisez l'historique du meilleur cours acheteur~: à chaque changement, sa nouvelle valeur, ou rien si le carnet est vide. Il doit s'exécuter en $O(n \log n)$. Que produisez-vous pour $(0, 5, 10)$, $(1, 3, 12)$, $(2, 4, 11)$ et $(6, 7, 9)$~?
\end{question}

\begin{question}[$\star\star\star\star$ \iqroles{QD, Dev, QR} \iqfirm{toutes} \iqlevel{confirmé}]\label{iq:iv3:algorithms:73}
Une liste de $n$ entiers peut contenir des négatifs. Comptez les segments contigus non vides dont la somme est comprise entre $L$ et $U$ (bornes incluses), en temps $O(n \log n)$. Quel résultat obtenez-vous pour \texttt{[3, -2, 1, 4]} avec $L = 1$ et $U = 3$~? Pourquoi une table de hachage des sommes préfixes antérieures ne suffit-elle pas, et pourquoi les pointeurs de votre étape de comptage ne font-ils qu'avancer~?
\end{question}

\begin{question}[$\star\star\star\star$ \iqroles{QR, QD, QT} \iqfirm{hedge fund, teneur de marché, banque} \iqlevel{confirmé}]\label{iq:iv3:algorithms:74}
Vous devez acheter $X$ actions sur $T$ périodes. L'achat de $q$ actions à la période $t$ coûte $a_t q + b_t q^2$, avec $b_t \ge 0$, et $q$ est un nombre entier. Trouvez la répartition la moins chère, en $O(X \log T)$. Pour $X = 5$, $a = (1, 3)$ et $b = (1, 1)$, quelle est la répartition et quel est son coût~? Démontrez que la méthode est juste, et dites où elle échoue si chaque période comporte aussi des frais fixes pour négocier.
\end{question}

\begin{question}[$\star\star\star\star$ \iqroles{QD, Dev, QR} \iqfirm{toutes} \iqlevel{confirmé}]\label{iq:iv3:algorithms:75}
Il y a des ordres d'achat et des ordres de vente, et une table indique quelles paires peuvent échanger, par exemple à cause de limites de crédit. Chaque ordre n'échange qu'avec au plus une contrepartie. Trouvez le plus grand nombre de paires en temps polynomial. L'acheteur 1 peut échanger avec les vendeurs A et B, et l'acheteur 2 seulement avec A. Que produit le fait d'apparier d'abord l'acheteur 1 à A, et comment votre méthode le corrige-t-elle~? Quel en est le coût~?
\end{question}

\begin{question}[$\star\star\star\star$ \iqroles{QD, Dev} \iqfirm{toutes} \iqlevel{confirmé}]\label{iq:iv3:algorithms:76}
Écrivez une fonction de correspondance de motifs où un point correspond à un caractère quelconque, et une étoile placée après un caractère (ou un point) correspond à zéro ou plusieurs de celui-ci, la chaîne entière devant correspondre. Faites-le en temps $O(mn)$ pour des longueurs de texte et de motif $n$ et $m$. \texttt{a*b*} correspond-il au texte \texttt{ba}~? Combien coûte un retour arrière naïf sur \texttt{a*a*a*a*b} appliqué à une longue suite de \texttt{a}~?
\end{question}

\begin{question}[$\star\star\star\star$ \iqroles{QD, Dev, QR} \iqfirm{banque, hedge fund} \iqlevel{confirmé}]\label{iq:iv3:algorithms:77}
Les contreparties ont des soldes nets dont la somme est nulle, un solde positif signifiant qu'on leur doit de l'argent. Un virement déplace un montant quelconque entre deux d'entre elles. Trouvez le plus petit nombre de virements qui ramène tous les soldes à zéro, pour au plus une vingtaine de contreparties. Quelle est la réponse pour \texttt{[5, -5, 3, -1, -2]}~? Pourquoi apparier de façon répétée le plus gros créancier au plus gros débiteur ne fonctionne-t-il pas toujours~?
\end{question}

\begin{question}[$\star\star\star\star$ \iqroles{QD, Dev, QR} \iqfirm{toutes} \iqlevel{confirmé}]\label{iq:iv3:algorithms:78}
Vous choisissez un sous-ensemble non vide de positions dans une liste de $n$ entiers, en ordre croissant, de sorte que deux positions choisies voisines soient distantes d'au plus $k$. Maximisez la somme des valeurs choisies en temps $O(n)$. Quel résultat obtenez-vous pour \texttt{[7, -2, -9, 4, 15]} avec $k = 2$~? Quelle est la récurrence, et quelle structure de données supprime son facteur $k$~?
\end{question}

\begin{question}[$\star\star\star\star$ \iqroles{QD, Dev} \iqfirm{toutes} \iqlevel{confirmé}]\label{iq:iv3:algorithms:79}
Étant donné une liste de $n$ entiers, renvoyez la somme, sur tous les segments contigus non vides, du minimum de chaque segment, en temps $O(n)$. Quel résultat obtenez-vous pour \texttt{[2, 4, 1, 3]}~? Qu'est-ce qui ne va pas avec les valeurs répétées, et comment comptez-vous chaque segment exactement une fois~? Essayez \texttt{[2, 2]}.
\end{question}

\begin{question}[$\star\star\star\star$ \iqroles{QR, QD} \iqfirm{hedge fund, teneur de marché} \iqlevel{confirmé}]\label{iq:iv3:algorithms:80}
On vous donne une liste de $n$ rendements quotidiens, ainsi qu'une longueur minimale $k$. Trouvez la plus grande moyenne sur des segments contigus d'au moins $k$ jours, en temps $O(n \log(R/\epsilon))$, où $R$ est l'étendue des valeurs et $\epsilon$ votre tolérance. Quelle est la réponse pour \texttt{[2, 10, -6, -5, 40, 4]} avec $k = 4$~? Pourquoi ne suffit-il pas de ne regarder que les segments d'exactement $k$ jours~?
\end{question}
